\chapter{Distributed Systems Fundamentals}
\label{ch:distributed-fundamentals}

\begin{chaptergoalsbox}
\textbf{What you will learn:}
\begin{itemize}[leftmargin=12pt]
  \item The inescapable physical limits of distributed computing: speed of light in fiber, cross-datacenter packet latency, jitter, and link saturation.
  \item Why networks are asynchronous and unreliable: packet loss, reordering, duplicate delivery, and network partitions.
  \item Time and order: physical clock drift, crystal oscillator mechanics, NTP vulnerabilities, Google TrueTime ($\varepsilon$-bounded uncertainty), Lamport logical timestamps, and Vector Clocks.
  \item Partial failure modes and theoretical models: synchronous vs. partially synchronous vs. asynchronous models, crash-stop vs. crash-recovery vs. Byzantine fault models.
\end{itemize}

\textbf{What you should already understand:}
\begin{itemize}[leftmargin=12pt]
  \item Basic networking protocols (TCP/IP handshake, UDP, round-trip time calculations).
  \item Single-node multithreading, concurrency hazards (race conditions, mutexes), and OS context switches.
\end{itemize}

\textbf{Why this topic exists:}
Single-node applications can rely on deterministic shared memory, synchronous bus communication, and monotonic hardware clocks. In distributed systems, none of these guarantees hold. Every assumption of shared state breaks down when computation crosses machine boundaries.

\textbf{Motivating practical problem:}
A multi-region e-commerce service experiences intermittent phantom inventory deductions and negative balances during Black Friday. Engineers discover that database replicas in Frankfurt and Virginia disagree on transaction order because they relied on system wall-clock timestamps (\texttt{System.currentTimeMillis()}) synchronized via standard public NTP servers drifting by up to 250 milliseconds.
\end{chaptergoalsbox}

\section{The Physical Reality of Distributed Computing}

Every distributed system is fundamentally constrained by physics. When software runs on a single multicore CPU, the distance between cores is measured in millimeters, and signals propagate across copper interconnects in fractions of a nanosecond. In contrast, when nodes communicate over a local area network (LAN) or wide area network (WAN), physical distance, optical refraction, and packet routing dictate strict lower bounds on latency.

\subsection{Speed of Light in Fiber and the Latency Floor}

The vacuum speed of light is $c \approx 299{,}792\text{ km/s}$. However, in standard single-mode optical fiber (such as Corning SMF-28), light travels through a silica glass core with a refractive index of approximately $n \approx 1.468$. The velocity of propagation $v$ in optical fiber is given by:

\begin{equation}
v = \frac{c}{n} \approx \frac{299{,}792\text{ km/s}}{1.468} \approx 204{,}218\text{ km/s} \approx 4.89\ \mu\text{s/km}
\end{equation}

For a transatlantic round trip between New York City and London (great-circle distance $\approx 5{,}570\text{ km}$, actual terrestrial/subsea fiber cable route $\approx 6{,}600\text{ km}$), the physical round-trip time (RTT) cannot be lower than:

\begin{equation}
\text{RTT}_{\text{min}} = \frac{2 \times 6{,}600\text{ km}}{204{,}218\text{ km/s}} \approx 64.6\text{ ms}
\end{equation}

In practice, packet switching, optical-electrical-optical (OEO) repeaters, Erbium-Doped Fiber Amplifiers (EDFAs), router queuing delays, and TCP stack processing elevate real-world transatlantic RTTs to between $70\text{ ms}$ and $85\text{ ms}$.

\begin{figure}[htbp]
\centering
\begin{tikzpicture}[
  scale=0.9,
  node distance=2.5cm,
  server/.style={rectangle, draw=headingblue, fill=codebg, thick, rounded corners=3pt, minimum width=2.8cm, minimum height=1.2cm, align=center, font=\sffamily\small},
  timeline/.style={thick, dashed, color=gray!60},
  packet/.style={-{Stealth[scale=1.0]}, thick, color=linkblue}
]
  \node[server] (client) at (0, 0) {\textbf{Client (US-East)}\\Virginia DC};
  \node[server] (server) at (8, 0) {\textbf{Replica (EU-West)}\\Frankfurt DC};

  \draw[timeline] (0, -0.8) -- (0, -6.5);
  \draw[timeline] (8, -0.8) -- (8, -6.5);

  \draw[packet] (0, -1.5) -- node[above, sloped, font=\sffamily\footnotesize] {SYN ($t=0\text{ ms}$)} (8, -2.5);
  \draw[packet] (8, -2.5) -- node[above, sloped, font=\sffamily\footnotesize] {SYN-ACK ($t=40\text{ ms}$)} (0, -3.5);
  \draw[packet] (0, -3.5) -- node[above, sloped, font=\sffamily\footnotesize] {ACK + TLS ClientHello ($t=80\text{ ms}$)} (8, -4.5);
  \draw[packet] (8, -4.5) -- node[above, sloped, font=\sffamily\footnotesize] {TLS ServerHello + App Data ($t=120\text{ ms}$)} (0, -5.5);

  \node[anchor=west, font=\sffamily\footnotesize, color=warningRed] at (8.5, -3.5) {$\Delta t_{\text{RTT}} \approx 80\text{ ms}$};
  \node[anchor=west, font=\sffamily\footnotesize, color=textdark] at (8.5, -5.5) {Minimum 2 RTTs before 1st byte: $160\text{ ms}$};
\end{tikzpicture}
\caption{Cross-Datacenter TCP/TLS Handshake Latency Budget across Subsea Fiber}
\label{fig:cross-dc-latency}
\end{figure}

\begin{intuitionbox}[The Asymmetry of Local vs. Remote State]
On a single machine, reading from L1 cache takes $\sim 0.5\text{ ns}$, reading from main memory (DRAM) takes $\sim 100\text{ ns}$, and reading from an NVMe SSD takes $\sim 10\text{ to }50\ \mu\text{s}$. A cross-datacenter network packet takes $80\text{ ms} = 80{,}000{,}000\text{ ns}$. In relative terms, if an L1 cache reference took 1 second, a round-trip network packet across the Atlantic would take over 5 years.
\end{intuitionbox}

\subsection{The Fallacies of Distributed Computing}

In 1994, L. Peter Deutsch and James Gosling articulated the ``Eight Fallacies of Distributed Computing''. Every engineer designing distributed infrastructure must internalize these mistaken assumptions:
\begin{enumerate}[leftmargin=18pt]
  \item \textbf{The network is reliable}: Networks drop packets due to buffer overruns, fiber cuts, faulty transceivers, BGP route flaps, and switch reboots.
  \item \textbf{Latency is zero}: Network hops always incur physical transit and queuing delays.
  \item \textbf{Bandwidth is infinite}: While backbone links achieve Terabits per second, intra-rack switches and egress transit pipes saturate under burst traffic, causing tail latency spikes.
  \item \textbf{The network is secure}: Distributed packets transit multi-tenant infrastructure, middleboxes, and public backbones; encryption and integrity verification (TLS/mTLS) are mandatory.
  \item \textbf{Topology doesn't change}: Servers scale horizontally, IP addresses migrate, network interfaces fail, and cloud providers reroute virtual networks dynamically.
  \item \textbf{There is one administrator}: Heterogeneous systems span multiple administrative domains, authorization tiers, and cloud availability zones with divergent policies.
  \item \textbf{Transport cost is zero}: Serialization, deserialization, kernel-to-user-space context switching, memory copies, and cloud egress bandwidth charges represent non-trivial compute and monetary costs.
  \item \textbf{The network is homogeneous}: Workloads execute across distinct hardware generations, CPU architectures, operating system kernels, and NIC drivers.
\end{enumerate}

\section{Time, Clocks, and Ordering}

Single-node programs rely on monotonically increasing hardware counters or system clocks to order events. When an application writes record $A$ at $t_1$ and record $B$ at $t_2$, the condition $t_1 < t_2$ guarantees that $A$ occurred before $B$. In distributed systems, this assumption is invalid.

\subsection{Physical Clocks and Clock Drift}

Computer motherboards track time using quartz crystal oscillators. Quartz crystals vibrate at specific resonant frequencies when subjected to an electric field (the piezoelectric effect). However, manufacturing variations, ambient temperature fluctuations, and crystal aging cause drift.

A standard server quartz oscillator drifts by approximately 1 to 50 parts per million (ppm). A drift of 10 ppm accumulates:
\begin{equation}
10\ \mu\text{s/s} = 0.864\text{ seconds per day} \approx 26\text{ seconds per month}
\end{equation}

\subsection{Network Time Protocol (NTP) and Its Limits}

To mitigate drift, machines periodically synchronize their clocks with reference time sources using the Network Time Protocol (NTP). In NTP, a client exchanges UDP packets with a stratum-1 or stratum-2 time server.

Let a client send a request at client time $t_0$, received by the server at server time $t_1$. The server transmits its response at server time $t_2$, which reaches the client at client time $t_3$.

\begin{figure}[htbp]
\centering
\begin{tikzpicture}[
  scale=0.9,
  timeline/.style={thick, color=headingblue},
  point/.style={circle, fill=linkblue, inner sep=1.5pt},
  packet/.style={-{Stealth[scale=1.0]}, thick, color=linkblue}
]
  \draw[timeline] (0, 0) -- (0, -5) node[below, font=\sffamily\small] {Client Timeline};
  \draw[timeline] (8, 0) -- (8, -5) node[below, font=\sffamily\small] {Server Timeline};

  \node[point] at (0, -1) {};
  \node[anchor=east, font=\sffamily\footnotesize] at (-0.2, -1) {Transmit: $t_0$};

  \node[point] at (8, -2) {};
  \node[anchor=west, font=\sffamily\footnotesize] at (8.2, -2) {Receive: $t_1$};

  \node[point] at (8, -3) {};
  \node[anchor=west, font=\sffamily\footnotesize] at (8.2, -3) {Transmit: $t_2$};

  \node[point] at (0, -4) {};
  \node[anchor=east, font=\sffamily\footnotesize] at (-0.2, -4) {Receive: $t_3$};

  \draw[packet] (0, -1) -- (8, -2);
  \draw[packet] (8, -3) -- (0, -4);

  \node[font=\sffamily\footnotesize, color=hoodSlate] at (4, -1.2) {Forward Delay: $d_1 = (t_1 - t_0) - \theta$};
  \node[font=\sffamily\footnotesize, color=hoodSlate] at (4, -3.2) {Reverse Delay: $d_2 = (t_3 - t_2) + \theta$};
\end{tikzpicture}
\caption{NTP Round-Trip Message Exchange and Clock Offset Calculation}
\label{fig:ntp-offset}
\end{figure}

The total round-trip network delay $\delta$ and clock offset $\theta$ are calculated as:
\begin{equation}
\delta = (t_3 - t_0) - (t_2 - t_1)
\end{equation}
\begin{equation}
\theta = \frac{(t_1 - t_0) + (t_2 - t_3)}{2}
\end{equation}

\begin{productionwarningbox}[The Asymmetry Fatal Flaw of NTP]
The offset formula strictly assumes that outbound network transit delay equals inbound transit delay ($d_1 = d_2$). If outbound network routing diverges from inbound routing (asymmetric BGP paths) or if outbound queuing differs from inbound queuing by $30\text{ ms}$, NTP will calculate an offset that is inaccurate by $15\text{ ms}$. Over the public Internet, NTP error commonly reaches $50\text{ ms}$ to $250\text{ ms}$. In virtualized cloud environments, VM hypervisor pauses, live migrations, and CPU throttling introduce artificial time dilation jumps of hundreds of milliseconds.
\end{productionwarningbox}

\subsection{Google TrueTime: Bounded Uncertainty}

To build Google Spanner---a globally distributed database offering external consistency without lock contention on read transactions---Google designed the TrueTime API.

Instead of exposing time as a discrete value ($t = \text{now}()$), TrueTime represents time as an interval $[t_{\text{earliest}}, t_{\text{latest}}]$ with guaranteed error bound $\varepsilon$:
\begin{equation}
\text{TrueTime.now}() \to [t_{\text{now}} - \varepsilon, t_{\text{now}} + \varepsilon]
\end{equation}

Google limits $\varepsilon$ to a small bound (typically $\le 7\text{ ms}$) by outfitting each datacenter with a dedicated time master cluster equipped with redundant GPS receivers and atomic rubidium clocks. If GPS fails or satellite signals are spoofed, the atomic clocks provide low-drift holdover timekeeping.

To ensure that transaction $T_1$ commit timestamp $s_1$ is strictly less than transaction $T_2$ commit timestamp $s_2$ whenever $T_2$ begins after $T_1$ commits, Spanner employs \textbf{commit wait}:

\begin{underthehoodbox}[Spanner's Commit Wait Rule]
When transaction $T_1$ prepares to commit, the coordinator picks absolute timestamp $s_1 = \text{TrueTime.now}().\text{latest}$. The coordinator refuses to release write locks or expose $T_1$'s mutated records to any client until it has paused for at least $2\varepsilon$:
\begin{equation}
\text{wait\_until}(\text{TrueTime.now}().\text{earliest} > s_1)
\end{equation}
This physical delay guarantees that any subsequent transaction $T_2$ will read $\text{TrueTime.now}().\text{earliest} > s_1$, mathematically enforcing linearizable ordering across global datacenters without cross-region coordination during read transactions.
\end{underthehoodbox}

\subsection{Lamport Logical Timestamps}

In 1978, Leslie Lamport published the seminal paper \emph{Time, Clocks, and the Ordering of Events in a Distributed System}, establishing that absolute physical time is unnecessary to order causally related events.

Lamport defined the \textbf{happens-before} relation ($\to$):
\begin{enumerate}[leftmargin=18pt]
  \item If event $a$ and event $b$ occur within the same process, and $a$ occurs before $b$, then $a \to b$.
  \item If event $a$ is the sending of a message by one process, and event $b$ is the receipt of that same message by another process, then $a \to b$.
  \item If $a \to b$ and $b \to c$, then $a \to c$ (transitivity).
\end{enumerate}

If neither $a \to b$ nor $b \to a$, the events are \textbf{concurrent} ($a \parallel b$).

\begin{lstlisting}[language=Python, caption={Implementation of a Lamport Logical Clock in Python}]
import threading

class LamportClock:
    def __init__(self, node_id: str):
        self.node_id = node_id
        self.counter = 0
        self._lock = threading.Lock()

    def tick(self) -> int:
        """Internal local event."""
        with self._lock:
            self.counter += 1
            return self.counter

    def send_event(self) -> int:
        """Invoked when dispatching a message over the network."""
        with self._lock:
            self.counter += 1
            return self.counter

    def receive_event(self, received_timestamp: int) -> int:
        """Invoked when receiving a message containing a remote timestamp."""
        with self._lock:
            self.counter = max(self.counter, received_timestamp) + 1
            return self.counter

    def compare(self, ts_a: int, node_a: str, ts_b: int, node_b: str) -> int:
        """Total order tie-breaking using unique node identifiers."""
        if ts_a < ts_b:
            return -1
        elif ts_a > ts_b:
            return 1
        else:
            return -1 if node_a < node_b else (1 if node_a > node_b else 0)
\end{lstlisting}

\begin{commonmistakebox}[Lamport Clocks Cannot Detect Concurrency]
A critical limitation of Lamport clocks is that while $a \to b \implies L(a) < L(b)$, the reverse is \textbf{not} true:
\begin{equation}
L(a) < L(b) \;\not\Longrightarrow\; a \to b
\end{equation}
If $L(a) = 3$ and $L(b) = 5$, it is impossible to determine whether $a$ caused $b$ or if $a$ and $b$ were completely concurrent events on isolated nodes. To detect and preserve concurrency, systems must use Vector Clocks.
\end{commonmistakebox}

\subsection{Vector Clocks}

A Vector Clock for a system of $N$ nodes is an array of $N$ logical integers: $V = [v_1, v_2, \dots, v_N]$. Node $i$ maintains its local vector $V_i$.

\textbf{Vector Clock Rules:}
\begin{enumerate}[leftmargin=18pt]
  \item Initially, all clock entries are zero: $V_i[j] = 0$ for all $j \in \{1, \dots, N\}$.
  \item Before node $i$ generates an internal or outbound event, it increments its own component:
  \begin{equation}
  V_i[i] \leftarrow V_i[i] + 1
  \end{equation}
  \item When node $i$ transmits a message, it attaches its current vector $V_i$.
  \item When node $j$ receives a message containing vector $V_{\text{msg}}$, it updates its clock:
  \begin{equation}
  V_j[k] \leftarrow \max(V_j[k], V_{\text{msg}}[k]) \quad \forall k \in \{1, \dots, N\}
  \end{equation}
  \begin{equation}
  V_j[j] \leftarrow V_j[j] + 1
  \end{equation}
\end{enumerate}

\textbf{Comparison and Causality Detection:}
\begin{itemize}[leftmargin=18pt]
  \item $V_A \le V_B \iff \forall k: V_A[k] \le V_B[k]$
  \item $V_A < V_B \iff (V_A \le V_B) \land (\exists k: V_A[k] < V_B[k])$
  \item If $V_A < V_B$, then $A$ causally preceded $B$ ($A \to B$).
  \item If neither $V_A \le V_B$ nor $V_B \le V_A$, then $A$ and $B$ are concurrent ($A \parallel B$). This indicates a conflict that must be resolved via application business logic (e.g., shopping cart merge) or Conflict-Free Replicated Data Types (CRDTs).
\end{itemize}

\begin{lstlisting}[language=Python, caption={Vector Clock Implementation with Conflict Detection in Python}]
from typing import Dict, Tuple

class VectorClock:
    def __init__(self, node_id: str):
        self.node_id = node_id
        self.clock: Dict[str, int] = {node_id: 0}

    def increment(self) -> Dict[str, int]:
        self.clock[self.node_id] = self.clock.get(self.node_id, 0) + 1
        return dict(self.clock)

    def update(self, remote_clock: Dict[str, int]) -> Dict[str, int]:
        # Merge by taking pairwise max
        for node, ts in remote_clock.items():
            self.clock[node] = max(self.clock.get(node, 0), ts)
        self.clock[self.node_id] = self.clock.get(self.node_id, 0) + 1
        return dict(self.clock)

    @staticmethod
    def compare(v1: Dict[str, int], v2: Dict[str, int]) -> str:
        """Returns: 'EQUAL', 'BEFORE', 'AFTER', or 'CONCURRENT'."""
        all_nodes = set(v1.keys()) | set(v2.keys())
        v1_greater = False
        v2_greater = False

        for node in all_nodes:
            val1 = v1.get(node, 0)
            val2 = v2.get(node, 0)
            if val1 > val2:
                v1_greater = True
            elif val2 > val1:
                v2_greater = True

        if v1_greater and not v2_greater:
            return "AFTER"       # v1 happened after v2
        elif v2_greater and not v1_greater:
            return "BEFORE"      # v1 happened before v2
        elif not v1_greater and not v2_greater:
            return "EQUAL"
        else:
            return "CONCURRENT"  # Conflict: both evolved independently!
\end{lstlisting}

\begin{realworldbox}[Amazon Dynamo and Riak Siblings]
Amazon DynamoDB's predecessor (the original 2007 Dynamo paper) and basho Riak utilized vector clocks to handle multi-master shopping cart writes during network partitions. When two datacenters were partitioned, a customer adding an item in California created vector $[CA:1, VA:0]$, while an item added in Virginia created $[CA:0, VA:1]$. Upon network recovery, comparing the two vectors returned \texttt{CONCURRENT}. Instead of discarding one write, the system preserved both entries as ``siblings'' and prompted the client application to union the items.
\end{realworldbox}

\section{Partial Failures and Theoretical Models}

In single-node software, components either execute or the operating system terminates the process (fail-stop). In distributed systems, failure is \textbf{partial}: some nodes continue operating normally while others crash, experience intermittent GC pauses, or become unreachable across broken network segments.

\subsection{Timing Models}

Theoretical distributed algorithms are proven under one of three fundamental timing assumptions:

\begin{table}[htbp]
\centering\small
\begin{tabular}{lp{4.2cm}p{4.2cm}p{3.8cm}}
\toprule
\textbf{Timing Model} & \textbf{Message Latency} & \textbf{Clock Drift} & \textbf{Real-World Feasibility} \\
\midrule
\textbf{Synchronous} & Strictly bounded: delay $\le D_{\text{max}}$ & Strictly bounded: drift $\le \rho$ & Specialized aerospace, avionics buses, hard real-time hardware. \\
\textbf{Asynchronous} & Unbounded: messages arrive eventually, no time limit & Clocks do not exist or drift arbitrarily & Most conservative theoretical model (FLM theorem applies). \\
\textbf{Partially Synchronous} & Asynchronous initially, but becomes bounded after an unknown Global Stabilization Time ($GST$) & Clocks drift, but drift becomes bounded after $GST$ & \textbf{Standard internet and cloud reality}; basis for Raft and Paxos. \\
\bottomrule
\end{tabular}
\caption{Distributed Systems Timing Models Comparison}
\label{tab:timing-models}
\end{table}

\subsection{Fault Models: Crash-Stop to Byzantine}

Node failure modes determine the complexity of consensus algorithms:

\begin{enumerate}[leftmargin=18pt]
  \item \textbf{Crash-Stop (Fail-Stop)}: A node operates correctly until it halts. Once halted, it never resumes or emits messages. Simplest model to reason about.
  \item \textbf{Crash-Recovery (Fail-Recovery)}: A node can halt at any time (e.g., kernel panic, power outage) and later reboot. When recovering, it retains state preserved on non-volatile storage (WAL, disk) but loses volatile RAM state. Raft and Paxos assume this model.
  \item \textbf{Omission Faults}: A node fails to transmit or receive messages due to full NIC buffers, dropped packets, or kernel socket exhaustion.
  \item \textbf{Byzantine Faults (Arbitrary Failures)}: Nodes may crash, drop messages, delay responses, forge cryptographic signatures, collude, or transmit intentionally conflicting data to different peers (split-brain sabotage). Common in blockchain and untrusted multi-party federations.
\end{enumerate}

\begin{figure}[htbp]
\centering
\begin{tikzpicture}[
  scale=0.85,
  box/.style={rectangle, draw=headingblue, fill=codebg, thick, rounded corners=3pt, minimum width=2.8cm, minimum height=1.1cm, align=center, font=\sffamily\footnotesize}
]
  \node[box, fill=exampleBg] (crash) at (0, 0) {\textbf{Crash-Stop}\\Clean termination};
  \node[box, fill=intuitionBg] (recovery) at (4, 0) {\textbf{Crash-Recovery}\\Reboot + Disk state};
  \node[box, fill=mistakeBg] (omission) at (8, 0) {\textbf{Omission}\\Dropped packets/NIC};
  \node[box, fill=warningBg] (byzantine) at (12, 0) {\textbf{Byzantine}\\Malicious / Corrupt};

  \draw[-{Stealth[scale=1.2]}, very thick, color=warningRed] (-1.5, -1.2) -- (13.5, -1.2)
    node[midway, below=3pt, font=\sffamily\small\bfseries] {Fault Model Severity and Algorithmic Complexity Spectrum $\longrightarrow$};
\end{tikzpicture}
\caption{Spectrum of Distributed Fault Models}
\label{fig:fault-spectrum}
\end{figure}

\begin{underthehoodbox}[The FLP Impossibility Result (Fischer, Lynch, Paterson 1985)]
The FLP theorem proves that in a fully \textbf{asynchronous} network model, no deterministic consensus protocol can guarantee both \textbf{Safety} (nothing bad happens, all nodes agree on the same value) and \textbf{Liveness} (something good happens eventually, nodes terminate and reach agreement) in the presence of even a \textbf{single unannounced crash-stop failure}.
Because it is impossible to distinguish between a crashed node and an arbitrarily slow node in an asynchronous model, protocols like Paxos and Raft trade off strict theoretical liveness: they ensure safety unconditionally under all network delays, but guarantee progress (liveness) only when the network transitions into a \textbf{partially synchronous} state.
\end{underthehoodbox}

\section{Summary \& Key Takeaways}
\begin{itemize}[leftmargin=18pt]
  \item Physical transmission speeds through optical fiber impose a strict latency floor of $\sim 4.89\ \mu\text{s/km}$, making zero-latency cross-region communication physically impossible.
  \item Physical wall clocks drift continuously due to thermal and mechanical oscillator variations. Public NTP cannot guarantee sub-millisecond precision due to asymmetric network delays.
  \item Google TrueTime solves absolute ordering by exposing an uncertainty interval $[t_{\text{now}} - \varepsilon, t_{\text{now}} + \varepsilon]$ and enforcing commit wait periods ($2\varepsilon$) backed by GPS receivers and atomic clocks.
  \item Lamport timestamps capture causality ($a \to b \implies L(a) < L(b)$) but cannot differentiate between causal and concurrent events.
  \item Vector clocks track independent node histories, allowing distributed systems to detect concurrent writes and data divergence.
  \item Real-world clouds operate under the \textbf{partially synchronous, crash-recovery} model. The FLP impossibility theorem proves that deterministic consensus in asynchronous networks must sacrifice unconditional liveness during network partitions.
\end{itemize}

\section{Exercises}

\subsection{Mechanics}
Given three nodes $A, B, C$ with initial vector clocks $[0, 0, 0]$:
\begin{enumerate}[leftmargin=18pt]
  \item Node $A$ executes an internal event $e_1$.
  \item Node $A$ sends message $m_1$ to Node $B$.
  \item Node $C$ executes an internal event $e_2$ independently.
  \item Node $B$ receives message $m_1$, then sends message $m_2$ to Node $C$.
  \item Node $C$ receives message $m_2$.
\end{enumerate}
Calculate the vector clock state at each event and determine whether event $e_2$ and the receipt of $m_1$ at Node $B$ are causal or concurrent.

\subsection{Debugging}
A financial audit log application records transactions using machine timestamps:
\begin{lstlisting}[language=Python]
record = {
    "account_id": acc_id,
    "delta": amount,
    "timestamp_ms": int(time.time() * 1000)
}
db.insert(record)
\end{lstlisting}
During a leap second insertion or hypervisor NTP step-adjustment, two debits of \$5,000 occur. The database deduplication engine uses \texttt{(account\_id, timestamp\_ms)} as a unique composite primary key. Explain the failure mode and provide a robust architectural fix.

\subsection{Architecture}
Design a distributed sequence ID generator (similar to Twitter Snowflake) capable of generating 100,000 globally unique, time-ordered 64-bit integers per second across 1,024 independent worker nodes without inter-node coordination during ID generation. Address clock backwards skew (NTP backward steps) explicitly.

\section{Detailed Solutions}

\subsection{Solution to Mechanics Exercise}
\begin{enumerate}[leftmargin=18pt]
  \item Node $A$ internal event $e_1$: Node $A$ increments its own index: $V_A = [1, 0, 0]$.
  \item Node $A$ sends $m_1$: increments $V_A = [2, 0, 0]$. Message carries vector $[2, 0, 0]$.
  \item Node $C$ internal event $e_2$: $V_C = [0, 0, 1]$.
  \item Node $B$ receives $m_1$: $V_B$ updates by taking max with $[2, 0, 0]$ and incrementing $B$'s slot: $V_B = [\max(0, 2), \max(0, 0) + 1, \max(0, 0)] = [2, 1, 0]$.
  \item Node $B$ sends $m_2$ to $C$: increments $V_B = [2, 2, 0]$. Message $m_2$ carries $[2, 2, 0]$.
  \item Node $C$ receives $m_2$: takes max with its current $[0, 0, 1]$ and increments $C$'s slot:
  $V_C = [\max(0, 2), \max(0, 2), \max(1, 0) + 1] = [2, 2, 2]$.
\end{enumerate}
\textbf{Causality evaluation:} Event $e_2$ has vector $[0, 0, 1]$. The receipt of $m_1$ at Node $B$ has vector $[2, 1, 0]$. Comparing the two: Node $A$'s and $B$'s components are greater in $B$'s vector, but Node $C$'s component is greater in $e_2$'s vector ($1 > 0$). Because neither vector dominates the other, $e_2 \parallel \text{receipt}(m_1)$ (\textbf{they are concurrent}).

\subsection{Solution to Debugging Exercise}
\textbf{Failure mode:} Physical wall clocks are not monotonic. When NTP steps backward (e.g., during clock slew adjustments or negative leap seconds), \texttt{time.time()} can return identical or decreasing values. If two debits occur during this window, the second debit will trigger a composite primary key collision (\texttt{UniqueViolationException}) and be rejected, or overwrite the prior record if an \texttt{UPSERT} pattern is used.
\textbf{Architectural fix:}
1. Never use wall-clock time as a primary key or unique identifier.
2. In single processes, use monotonic hardware timers (\texttt{time.monotonic\_ns()} in Python, \texttt{CLOCK\_MONOTONIC} in Linux).
3. Combine monotonic time with a node identifier and an auto-incrementing atomic sequence counter (e.g., Snowflake ID format: 41 bits timestamp, 10 bits worker ID, 12 bits sequence number).
4. Configure local NTP daemons with \textbf{clock slewing only} (passing \texttt{-x} to \texttt{ntpd} or configuring \texttt{chrony} with \texttt{makestep 0}) to forbid hard backward step jumps.

\subsection{Solution to Architecture Exercise}
\textbf{Bit Allocation Layout (64-bit integer):}
\begin{itemize}[leftmargin=18pt]
  \item \textbf{1 bit sign}: Always 0 (keeps integers positive).
  \item \textbf{41 bits millisecond timestamp}: Difference from custom epoch (e.g., Jan 1, 2026). Yields $(2^{41} - 1) / (1000 \times 3600 \times 24 \times 365) \approx 69.7$ years of lifespan.
  \item \textbf{10 bits worker/machine ID}: Supports up to $2^{10} = 1{,}024$ distinct worker instances. Assigned via ZooKeeper, Consul, or K8s StatefulSet pod ordinal index.
  \item \textbf{12 bits sequence counter}: Supports $2^{12} = 4{,}096$ IDs per millisecond per node ($4{,}096{,}000$ IDs/second per worker, well beyond the 100k requirement).
\end{itemize}

\textbf{Handling Clock Backward Skew:}
\begin{lstlisting}[language=Python]
class SnowflakeGenerator:
    def __init__(self, worker_id: int, epoch: int = 1767225600000):
        assert 0 <= worker_id < 1024
        self.worker_id = worker_id
        self.epoch = epoch
        self.sequence = 0
        self.last_timestamp = -1

    def next_id(self) -> int:
        current_time = self._current_ms()

        # Handle backward clock drift
        if current_time < self.last_timestamp:
            drift = self.last_timestamp - current_time
            if drift < 5:  # Tolerate small drift (<5ms) by spinning
                while current_time < self.last_timestamp:
                    time.sleep(0.001)
                    current_time = self._current_ms()
            else:
                raise SystemClockException(f"Clock moved backwards by {drift}ms")

        if current_time == self.last_timestamp:
            self.sequence = (self.sequence + 1) & 0xFFF
            if self.sequence == 0:
                # Sequence exhausted for this millisecond; spin wait
                while current_time <= self.last_timestamp:
                    current_time = self._current_ms()
        else:
            self.sequence = 0

        self.last_timestamp = current_time
        return ((current_time - self.epoch) << 22) | (self.worker_id << 12) | self.sequence

    def _current_ms(self) -> int:
        return int(time.time() * 1000)
\end{lstlisting}

\section{Further Reading \& Citations}
\begin{itemize}[leftmargin=18pt]
  \item Lamport, L. (1978). \emph{Time, clocks, and the ordering of events in a distributed system}. Communications of the ACM, 21(7), 558--565.
  \item Corbett, J. C., et al. (2013). \emph{Spanner: Google's globally distributed database}. ACM Transactions on Computer Systems (TOCS), 31(3), 1--22.
  \item Fischer, M. J., Lynch, N. A., \& Paterson, M. S. (1985). \emph{Impossibility of distributed consensus with one faulty process}. Journal of the ACM (JACM), 32(2), 374--382.
  \item Deutsch, L. P. (1994). \emph{The Eight Fallacies of Distributed Computing}. Sun Microsystems Technical Report.
\end{itemize}
